\documentclass[11pt]{article}
\usepackage[spanish]{babel}
\usepackage{graphicx}
\usepackage[margin=2cm]{geometry}
\usepackage{csquotes}
\usepackage[backend=biber,style=ieee]{biblatex}
\usepackage{hyperref}
\graphicspath{{figuras/}}
\addbibresource{referencias.bib}
\title{Notas de la exposición: computación paralela}
\author{GIRALDO CAGUA, HERRERA BARCO\\GARCIA SANGOQUIZA, LOMBEIDA ESCALERAS}
\date{}
\begin{document}
\maketitle
Tema: computación paralela, con ejemplos generales y aplicaciones a los PFC
de los integrantes.\par
\paragraph{Declaración de uso de IA.}
\textbf{Bryan Humberto Lombeida Escaleras:} utilicé OpenAI GPT-6.2 Sol para
buscar y contrastar fuentes, preparar y revisar mis diapositivas y notas,
y generar y revisar código, pruebas y documentación de la práctica en Java.
Las capturas documentan las ejecuciones de la práctica con datos de prueba;
el fallo se simuló y no corresponde a un incidente de SGRA.
La declaración general del equipo queda por completar con los aportes de los
demás integrantes.
\section*{GIRALDO CAGUA — Presentación del sistema (1 min) y arquitectura (2 min)}
[PENDIENTE: guion sustentado de tres minutos.]

\section*{HERRERA BARCO — Transparencias (1,5 min) y modelos de comunicación (1,5 min)}
[PENDIENTE: guion sustentado de tres minutos.]

\section*{GARCIA SANGOQUIZA — CAP y PACELC (2 min) y comparación crítica (1 min)}

\subsection*{CAP: ejemplo general y aplicación propuesta a CLAB (unos 75 s)}
Ejecutar tareas en paralelo no elimina la necesidad de coordinar los datos
compartidos. CAP estudia un límite de esa coordinación distribuida: durante
una partición de red no podemos garantizar simultáneamente consistencia y
disponibilidad~\cite{gsGilbertLynch2002,gsIbmCap}.
C significa que las operaciones se comportan como sobre una única copia
actualizada; A, que toda petición a un nodo operativo termina con una respuesta
válida, no simplemente con un error; P contempla que los nodos pierdan comunicación.

Imaginemos dos réplicas que venden la última unidad. Si quedan aisladas y ambas
confirman la venta, pueden venderla dos veces. Para evitarlo, alguna debe dejar
de confirmar mientras no pueda coordinarse.

En CLAB, propongo el mismo razonamiento para dos nodos que intentan reservar
el mismo laboratorio y horario: el nodo aislado no confirma una reserva sin
coordinarla. Se prioriza consistencia sobre disponibilidad de esa operación.
Es una propuesta, no una función implementada: la base E1 solo consulta
Laboratorios desde Reservas por HTTP; todavía no tiene reservas ni
persistencia~\cite{gsClabE1}. Dos servicios no demuestran paralelismo ni permiten
clasificar automáticamente todo CLAB como CP.

\subsection*{PACELC: consistencia frente a latencia (unos 45 s)}
PACELC complementa CAP: con partición, disponibilidad frente a consistencia;
sin partición, latencia frente a consistencia en datos
replicados~\cite{gsAbadi2012}.
Latencia es cuánto tarda una operación en responder. Por ejemplo, leer un precio
en una réplica cercana puede reducir la espera, pero devolver un valor antiguo.
Garantizar una lectura actualizada puede exigir comunicación adicional.

Para CLAB propongo consultas paralelas de disponibilidad sobre réplicas:
servirían como orientación, no como confirmación. Aunque una consulta muestre
libre un laboratorio, la reserva debe validarse y confirmarse de forma atómica
en el servicio autorizado. Son decisiones propuestas por operación, no una
etiqueta PACELC del CLAB actual. Paralelizar consultas no elimina el coste de
coordinarlas~\cite{gsAbadi2012,gsMicrosoftConsistency2025}.

\subsection*{Comparación crítica (1 min)}
[PENDIENTE: contrastar dos enfoques tecnológicos y ajustar el guion completo
a tres minutos mediante ensayo.]

\paragraph{Declaración de uso de IA de esta sección.}
Se utilizó OpenAI GPT-6.1 Sol para preparar este borrador de diapositivas y notas
de CAP y PACELC, consultar fuentes y revisar su correspondencia con el código
de CLAB.
La revisión personal, la adaptación del guion y el ensayo del expositor están
pendientes. Los ejemplos son ilustrativos, no resultados experimentales.

\section*{LOMBEIDA ESCALERAS — Tolerancia a fallos (2 min) y lecciones/conclusiones (1 min)}
[PENDIENTE: guion sustentado de tres minutos.]
\section*{Tabla de aportes}
\noindent\begin{tabular}{@{}p{0.27\linewidth}p{0.69\linewidth}@{}}
\textbf{Integrante} & \textbf{Trabajo en el grupo y participación en la exposición} \\
\hline
Bryan Humberto Lombeida Escaleras &
Organizó la plantilla Beamer y los archivos principales del equipo.
Preparó las diapositivas y notas 10--12, revisó fuentes y reunió evidencias
de lectura y de la práctica Java aplicada al flujo de SGRA.
Expone tolerancia a fallos, detección, recuperación y lecciones para el PFC. \\
\end{tabular}

\noindent\textit{Filas por completar: Giraldo Cagua, Herrera Barco y Garcia Sangoquiza.}
\printbibliography
\end{document}
